\documentclass{article}
\usepackage[margin=1in]{geometry}
\usepackage{cite}
\usepackage{amsmath,amssymb,amsfonts}
\usepackage{algorithmic}
\usepackage{graphicx}
\usepackage{textcomp}
\usepackage{xcolor}
\usepackage{soul}

\usepackage{multirow}
\usepackage{semantic}
\usepackage{array}

\usepackage{hyperref}
\hypersetup{
    colorlinks,
    linkcolor={red!50!black},
    citecolor={blue!50!black},
    urlcolor={blue!80!black}
}
\usepackage[capitalize, noabbrev, nameinlink]{cleveref}




\usepackage{macros}

\begin{document}


\newif\ifdraft
% For final version, replace next line by \draftfalse
%\drafttrue % for draft version
\drafttrue
\newcommand{\comment}[1]{\ifdraft{\textrm{\color{blue}[#1]}}\fi}
\newcommand{\john}[1]{\ifdraft{\textrm{\color{red}[#1 --John]}}\fi}

\title{GoosePEG
% \thanks{Identify applicable funding agency here. If none, delete this.}
}

\maketitle

\section{Syntax}

The syntax for GoosePEG expressions Expr is given in \cref{fig:syntax}.
We model nonterminals as a set of variables NonTerm, 
where the nonterminals for a particular grammar consist of a total function $F$ from nonterminal variables to expressions (i.e. 
$F : \text{NonTerm} \to \text{Expr}$).

\begin{figure}[h]
\centering

    \begin{tabular}{l c >{$}l<{$}}
        \multicolumn{3}{c}{c $\in$ Term, \quad v $\in$ NonTerm} \\
        $e, e_1, e_2 \in$ Expr & \bnfdef & \texttt{term} ~ c \\
                     & \bnfalt & \texttt{nterm} ~ v \\
                     & \bnfalt & e_1 \cdot e_2 \\
                     & \bnfalt & e_1 / e_2 \\
                     & \bnfalt & e^{*} \\
                     & \bnfalt & !e \\
    \end{tabular}
    %
    \caption{Core syntax for PEG grammars.}
    \label{fig:syntax}
\end{figure}

\section{Semantics}

The semantics of parsing expressions is given in \cref{fig:nat-semantics}.
Following Ford, we write $\parseto{e}{s}{r}$ to mean an expression $e$ parses a string $s$ to a result $r$, which can either be a string remainder $s'$ or a parse failure $\bot$.
For strings, we write $\epsilon$ for the empty string and list syntax otherwise (i.e., $c : t$ denotes a character $c$ and a tail $t$, and $s_1 \listcat s_2$ denotes the list concatenations of two strings $s_1$ and $s_2$).

Two rules. 

\begin{enumerate}
\item Normal rule just looks up from the memo table with precedence $n$: $\Gamma |- (e^{n}, s) \Rightarrow s'$
\item Updating rule changes the memo table:  $\Gamma |- (e^{n}, s) \to (s', \Gamma')$

\end{enumerate}

\begin{figure}
    \centering
    \[
    \begin{array}{c c c}
        \inference{}{\parseto{\texttt{term} ~ c}{c : s}{s}}{\textbf{t.n1}} & 
        \inference{c \neq c'}{\parseto{\texttt{term} ~ c}{c' : s}{\bot}}{\textbf{t.n2}} &
        \inference{}{\parseto{\texttt{term} ~ c}{[]}{\bot}}{\textbf{t.n3}} 
        \\ \\
        \inference{\parseto{F ~ v}{s}{r}}{\parseto{\texttt{nonterm} ~ v}{s}{r}}{\textbf{nt.n1}} &
        \inference{}{\parseto{\texttt{nonterm} ~ v}{s}{\bot}}{\textbf{nt.n2}} &
        \inference{
            \parseto{e}{s}{\bot} 
        }{\parseto{!e}{s}{s}}{\textbf{neg.n1}} 
        \\ \\
        \inference{
            \parseto{e}{s}{s'} 
        }{\parseto{!e}{s}{\bot}}{\textbf{neg.n2}} &
        \inference{
            \parseto{e_1}{s_1 \listcat s_2}{s_2} \\
            \parseto{e_2}{s_2}{r}
        }{\parseto{e_1 \cdot e_2}{s_1 \listcat s_2}{r}}{\textbf{cat.n1}} &
        \inference{
            \parseto{e_1}{s}{\bot} 
        }{\parseto{e_1 \cdot e_2}{s}{\bot}}{\textbf{cat.n2}} 
        \\ \\
        \inference{
            \parseto{e_1}{s}{s'} 
        }{\parseto{e_1 / e_2}{s}{s'}}{\textbf{alt.n1}} &
        \inference{
            \parseto{e_1}{s}{\bot} \\
            \parseto{e_2}{s}{r} 
        }{\parseto{e_1 / e_2}{s}{r}}{\textbf{alt.n2}} &
        \inference{
            \parseto{e}{s}{\bot} 
        }{\parseto{e^{*}}{s}{s}}{\textbf{rep.n1}} 
        \\ \\
        \inference{
            \parseto{e}{s}{s'} \\
            \parseto{e^{*}}{s'}{s''}
        }{\parseto{e^{*}}{s}{s''}}{\textbf{rep.n2}}
        & & \\ \\
    \end{array}
    \]
    %
    \caption{Core Natural semantics.}
    \label{fig:nat-semantics}
\end{figure}


\end{document}

